\parthead{IV}{Recovering the city around the approach}\label{part:4}
A current map cannot show a vanished building simply by being more detailed. For an approach in 1993, that absence matters: the object missing from today's image may be the object the aircraft had to pass. The historical problem begins with matching a place to a date.

Survey sheets offer one way in. The October 1989 and July 1994 editions of sheet 11-NE-11A place streets and mapped forms in the same Hong Kong 1980 Grid rectangle. Their labels give recognizable points of reference: Shing Tung House, the Regal Airport Hotel and Kowloon City Market. Viewed together, the editions allow individual forms and their surroundings to be compared across time \citep{maps}. A named form provides a lead; its presence, height and identity at an intervening date require additional evidence.

\fig{06_maps}{1}{Kowloon in survey sheet 11-NE-11A: the October 1989 field above, and matched 150 m-wide July 1994 windows below. Named locators let the reader compare mapped forms across editions. The sheets use the Hong Kong 1980 Grid; registered raster spacing is 0.25 m, not a positional-error bound. Matching forms are not automatically accepted as unchanged 1993 buildings. Source: Lands Department, HKSAR Government; CSDI Portal and DATA.GOV.HK, Old Survey Sheets. Government ownership retained; supplied as-is, without endorsement.}{fig:maps}

A photograph contributes a different kind of evidence. It records visible surfaces together at one exposure. A roof, a road junction and the pattern of adjoining blocks can be examined in their original setting, without first reducing them to polygons. The overlapping frames A36142 and A36143 are catalogued as 2 November 1993, two days before CAL605; A36143's apparent handwritten 3 November label remains an unresolved discrepancy \citep{photos}.

\subsection*{Depth from a displacement}
At Kowloon City Market, the roof feature marked R can be picked out in both photographs; S marks a rooftop service structure. Their displacement between the images provides depth information. Once the camera geometry and metric scale are fixed, the corresponding sight lines constrain each feature's position in three dimensions. In the reconstruction, R lies 10.77 m above the fitted local park/road plane and S lies 19.20 m above it. Both examples lie outside the map-control hull: their apparent detail exceeds the direct positional support available at their locations.

Here the image matches are anchored to manually interpreted ground features and survey-road controls. Camera and scan assumptions enter that alignment. The result retains 1,318 sparse matches and 171,307 filtered dense points. The number of points gives a sense of the recovered detail; the distribution of control determines how confidently that detail can be placed.

\fig{07_stereo}{1}{From historical photographs to recovered surface structure. The stereo reconstruction derives from overlapping A36142/A36143 images and map controls; the geographic interpretation is tied to the dated survey forms. Heights are relative to a fitted local park/road plane, not HKPD or height above each building's own ground. The control hull and missing geometry remain visible. The labelled roof features lie outside the control hull. Sources: Lands Department photographs and survey sheets, with the project's retained stereo/mesh reconstruction. Government/CSDI/DATA.GOV.HK attribution and ownership retained.}{fig:stereo}

Only 3,780 dense points lie inside the map-control hull, about 2.21\% of the retained cloud. Most of the reconstructed surface therefore extends beyond the small region directly enclosed by the alignment controls. The 2.21\% measures that support geometry; it is not an estimate of the fraction of the city reconstructed correctly. A detailed-looking surface can carry much less direct positional support than its appearance suggests.

The held-out horizontal control discrepancy is 5.754 m RMS. That number describes a horizontal registration diagnostic. It supplies no corresponding vertical roof-error bound. The vertical coordinate is conditional on a fitted local park/road plane, so a high point in the rendering means high relative to that plane. It is not automatically a building's height above its own base, still less an absolute roof elevation in a surveyed datum.

\FloatBarrier
A conservative mesh joins the supported geometry while leaving holes and boundaries; it contains 86,061 boundary edges. Those edges mark the limits of the reconstructed surface, including where views or correspondences run out. Closing a hole with a plausible wall would make the scene easier to recognize, but the new wall would come from the illustrator. It would add no observation. The retained mesh therefore remains non-watertight, and its detailed counts and earlier support correction remain in Appendix~\ref{app:sources}.

\subsection*{A height needs a vertical origin}
Some survey sheets contain approximate spot elevations above Hong Kong Principal Datum (HKPD). The settled obstacle inventory has 24 such readings at 12 sites on sheet 11-NW-10D, from November 1989 and June 1995. Repeated values can be carried forward between map editions; they are not necessarily independent measurements.

An elevation and a building height answer different questions. If a roof elevation and the underlying terrain already use HKPD, adding the roof number to the terrain number counts the vertical origin twice. Conversely, CAD's runway elevation of 15 ft AMSL, or 4.572 m in that stated reference, cannot simply be set equal to an HKPD elevation. A common vertical transformation still has to be justified. The project's present reconstruction has not accepted one.

Modern terrain and building products help locate potential evidence, but their dates remain attached to them. The DTM uses 2014--2015 photography and a 2019 revision; the building snapshot is from September 2026 \citep{dtm,buildings}. Neither supplies a complete census of what stood around the 1993 approach. Appendix~\ref{app:sources} preserves the full source inventory, dated matching counts and obstacle limitations.

The present historical geometry contains no accepted complete obstacle volumes or physical obstacle upper bounds. Its unknown areas remain marked in the exported scene. To test clearance, those missing pieces matter as much as the visible ones: a real aircraft sweeps wings, tail and landing gear through space, with tracking error around any chosen reference path. The reconstruction has recovered a dated fragment of the city and made its geometric support inspectable. Turning that fragment into a clearance model requires additional observations.

Even a fully recovered city would leave another moving part. The aircraft travels through air that need not move with the ground.
